\subsection*{Oblivious Transfer}
\label{sub:ObliviousTransfer}
Consider the following scenario:
\begin{quotation}
    Bob has a database with 10 entries.
    Alice would like to query it.
    \textit{Oblivious Transfer (OT)} allows Alice to retreive an entry of her choosing without Bob learning which one.
    Bob is guaranteed that Alice can only see the single entry she chose.
    She will not get any information about the other entries.
\end{quotation}

Yadav et al. \cite{yadav2022survey} describe \textit{Oblivious Transfer Protocols (OT)} as crucial tools for securely executing multi party computations.
It comes in three flavors: \textit{1-out-of-2 OT}, \textit{1-out-of-n OT}, and \textit{k-out-of-n OT}.
This research will focus on \textit{1-out-of-n PIR} as it can be used in \textit{Private Information Retreival} as described in the following subsection.
Notations and definitions are taken from Yadav et al. \cite{yadav2022survey}.

On a high level \textit{1-out-of-n OT} describes the data exchange of two parties, a sender and a receiver.
The sender has $n$ inputs $(x_0,...,x_{n-1})\in \{0,1\}^l$.
The receiver has a selection number $r\in [0,...,n-1]$.
After protocol execution the receiver learns $x_r$ while the sender learns nothing about $r$.
Note, that the receiver does not learn anything about $x_\psi$ with $\psi\in [0,...,n-1] \text{ and } \psi\neq r$.

\subsection*{Private Information Retreival (PIR)}
\label{sub:PrivateInformationRetreival}
The difference between oblivious transfer and \textit{Private Information Retreival (PIR)} \cite{gasarch2004survey} can be explained by taking another perspective on the scenario from above:
\begin{quotation}
    Again, assume Bob's database to have 10 entries.
    Alice would like to query it again.
    PIR allows Alice to retreive an element of her choosing.
    Bob will not learn which element Alice was interested in.
    Alice, however, has access to the full database.
\end{quotation}

The most trivial PIR variant is Alice retreiving the whole database \cite{gasarch2004survey}.
More advanced PIR-protocols fulfil this requirement in sublinear runtime.
\textit{DUORAM} (\ref{sec:3pDuoram}) uses DPFs to implement PIR and PIR-write operations in sub-linear runtime \cite{vadapalli2023duoram}.
